\Article{236}
\textsc{Problem.} \textit{Given two forms whose determinants are either equal or at least have the ratio of squares to one another, find a form composed of them.}

\textit{Solution.} Let the forms to be composed be $(a,b,c)\ldots f$, $(a',b',c')\ldots f'$; let their determinants be $d,d'$; let the greatest common divisors of the numbers $a,2b,c;\ a',2b',c'$ be $m,m'$, respectively; let the greatest common divisor of the numbers $dm'm',d'mm$, taken with the same sign as $d,d'$,
\SourceLeaf{254}
be $D$. Then $\frac{dm'm'}{D},\frac{d'mm}{D}$ are positive, relatively prime numbers whose product is a square; so they are themselves squares (Article 21). Thus $\sqrt{\frac{d}{D}},\sqrt{\frac{d'}{D}}$ are rational quantities, which we put $=n,n'$, taking for $n$ the positive or negative value according as the form $f$ is to enter the composition directly or inversely; similarly, we determine the sign of $n'$ from the way in which $f'$ is to enter the composition. Thus $mn',m'n$ are relatively prime integers; but $n$ and $n'$ may also be fractions. Having done this, we observe that $an',cn',a'n,c'n,bn'+b'n,bn'-b'n$ are integers. This is evident for the first four (since $an'=\frac{a}{m}mn'$, etc.); for the other two, it is proved in the same way as we proved in the preceding article that $R$ and $S$ are divisible by $e$.

Now take any four integers $\mathfrak{Q},\mathfrak{Q}',\mathfrak{Q}'',\mathfrak{Q}'''$, with only the condition that the four quantities on the left in the following equation (I) are not all $=0$ at once, and put
\[
\begin{aligned}
\mathfrak{Q}'an'+\mathfrak{Q}''a'n+\mathfrak{Q}'''(bn'+b'n)&=\mu q\qquad\ldotp\,\ldotp\,\ldotp\,\ldotp\,\ldotp\qquad\text{(I)}\\
-\mathfrak{Q}an'+\mathfrak{Q}'''c'n-\mathfrak{Q}''(bn'-b'n)&=\mu q'\\
\mathfrak{Q}'''cn'-\mathfrak{Q}a'n+\mathfrak{Q}'(bn'-b'n)&=\mu q''\\
-\mathfrak{Q}''cn'-\mathfrak{Q}'c'n-\mathfrak{Q}(bn'+b'n)&=\mu q'''
\end{aligned}
\]
so that $q,q',q'',q'''$ are integers with no common divisor. We achieve this by taking for $\mu$ the greatest common divisor of the four numbers on the left in these equations. Then, by Article 40, four integers $\mathfrak{P},\mathfrak{P}',\mathfrak{P}'',\mathfrak{P}'''$ can be found such that
\[
\mathfrak{P}q+\mathfrak{P}'q'+\mathfrak{P}''q''+\mathfrak{P}'''q'''=1
\]
Now determine the numbers $p,p',p'',p'''$ by the following equations:
\[
\begin{aligned}
\mathfrak{P}'an'+\mathfrak{P}''a'n+\mathfrak{P}'''(bn'+b'n)&=p\qquad\ldotp\,\ldotp\,\ldotp\,\ldotp\,\ldotp\qquad\text{(II)}\\
-\mathfrak{P}an'+\mathfrak{P}'''c'n-\mathfrak{P}''(bn'-b'n)&=p'\\
\mathfrak{P}'''cn'-\mathfrak{P}a'n+\mathfrak{P}'(bn'-b'n)&=p''\\
-\mathfrak{P}''cn'-\mathfrak{P}'c'n-\mathfrak{P}(bn'+b'n)&=p'''
\end{aligned}
\]
Finally, put
\[
\begin{gathered}q'q''-qq'''=Ann',\\pq'''+qp'''-p'q''-q'p''=2Bnn',\\p'p''-pp'''=Cnn'\end{gathered}
\]
\SourceLeaf{255}
Then $A,B,C$ are integers, and the form $(A,B,C)\ldots F$ is composed of the forms $f,f'$.

\textit{Proof.} I. From equations I, the following four equations are confirmed without difficulty:
\[
\begin{aligned}
0&=q'cn'-q''c'n-q'''(bn'-b'n)\qquad\ldotp\,\ldotp\,\ldotp\,\ldotp\,\ldotp\,\ldotp\qquad\text{(III)}\\
0&=qcn'+q'''a'n-q''(bn'+b'n)\\
0&=q'''an'+qc'n-q'(bn'+b'n)\\
0&=q''an'-q'a'n-q(bn'-b'n)
\end{aligned}
\]

II. Now suppose that the integers $\mathfrak{A},\mathfrak{B},\mathfrak{C},\mathfrak{A}',\mathfrak{B}',\mathfrak{C}',\mathfrak{N},\mathfrak{N}'$ have been determined so that
\[
\begin{aligned}
\mathfrak{A}a+2\mathfrak{B}b+\mathfrak{C}c&=m\\
\mathfrak{A}'a'+2\mathfrak{B}'b'+\mathfrak{C}'c'&=m'\\
\mathfrak{N}m'n+\mathfrak{N}'mn'&=1
\end{aligned}
\]
Then we have
\[
\begin{aligned}\mathfrak{A}a\mathfrak{N}'n'&+2\mathfrak{B}b\mathfrak{N}'n'+\mathfrak{C}c\mathfrak{N}'n'\\&+\mathfrak{A}'a'\mathfrak{N}n+2\mathfrak{B}'b'\mathfrak{N}n+\mathfrak{C}'c'\mathfrak{N}n=1\end{aligned}
\]
From this and equations (III), we readily confirm that if we put
\[
\begin{aligned}
-q'\mathfrak{A}\mathfrak{N}'-q''\mathfrak{A}'\mathfrak{N}-q'''(\mathfrak{B}\mathfrak{N}'+\mathfrak{B}'\mathfrak{N})&=\mathfrak{q}\\
q\mathfrak{A}\mathfrak{N}'-q'''\mathfrak{C}'\mathfrak{N}+q''(\mathfrak{B}\mathfrak{N}'-\mathfrak{B}'\mathfrak{N})&=\mathfrak{q}'\\
-q'''\mathfrak{C}\mathfrak{N}'+q\mathfrak{A}'\mathfrak{N}-q'(\mathfrak{B}\mathfrak{N}'-\mathfrak{B}'\mathfrak{N})&=\mathfrak{q}''\\
q''\mathfrak{C}\mathfrak{N}'+q'\mathfrak{C}'\mathfrak{N}+q(\mathfrak{B}\mathfrak{N}'+\mathfrak{B}'\mathfrak{N})&=\mathfrak{q}'''
\end{aligned}
\]
we have
\[
\begin{aligned}
\mathfrak{q}'an'+\mathfrak{q}''a'n+\mathfrak{q}'''(bn'+b'n)&=q\qquad\ldotp\,\ldotp\,\ldotp\,\ldotp\,\ldotp\,\ldotp\qquad\text{(IV)}\\
-\mathfrak{q}an'+\mathfrak{q}'''c'n-\mathfrak{q}''(bn'-b'n)&=q'\\
\mathfrak{q}'''cn'-\mathfrak{q}a'n+\mathfrak{q}'(bn'-b'n)&=q''\\
-\mathfrak{q}''cn'-\mathfrak{q}'c'n-\mathfrak{q}(bn'+b'n)&=q'''
\end{aligned}
\]

Whenever $\mu=1$, these equations are unnecessary; in their place we can keep equations (I), to which they are entirely analogous. If we now work out from equations II, IV the values of $Ann',2Bnn',Cnn'$ (\textit{i.e.}, of the numbers $q'q''-qq'''$, etc.), and remove the terms that cancel one another, we find that the terms of each
\SourceLeaf{256}
are products of integers either with $nn'$, with $dn'n'$, or with $d'nn$, and that all the terms making up $2Bnn'$ also contain the factor 2. From this (since $dn'n'=d'nn$, and so $\frac{dn'n'}{nn'}=\frac{d'nn}{nn'}=\sqrt{dd'}$ are integers) we conclude that $A,B,C$ are integers. \textit{Q. E. P.}

III. Substituting from equations (II) the values of $p,p',p'',p'''$, we readily verify, with the help of equations (III) and of
\[
\mathfrak{P}q+\mathfrak{P}'q'+\mathfrak{P}''q''+\mathfrak{P}'''q'''=1
\]
that
\[
\begin{aligned}
pq'-qp'&=an',\\ pq'''-qp'''-p'q''+q'p''&=2bn',\\p''q'''-q''p'''&=cn'\\
pq''-qp''&=a'n,\\ pq'''-qp'''+p'q''-q'p''&=2b'n,\\p'q'''-q'p'''&=c'n
\end{aligned}
\]
These equations are identical to the first six ($\Omega$) of the preceding article; the other three already hold by hypothesis. So (\textit{at the end of the same article}) the form $F$ passes into $ff'$ by the substitution $p,p',p'',p''';\ q,q',q''\,q'''$; its determinant is $=D$, or equal to the greatest common divisor of the numbers $dm'm',d'mm$, and therefore, by the fourth conclusion of the preceding article, $F$ is composed of $f,f'$. \textit{Q. E. S.} Finally, it is easy to see that $F$ is composed of $f,f'$ \textit{in the way} required, since the signs of the quantities $n,n'$ were determined correctly from the start.
